\section{Cost analysis}\label{sec:cost}

\subsection{Per-second cost comparison}

\begin{table}[htbp]
\centering
\caption{Per-substrate cost comparison.}
\label{tab:cost_per_substrate}
\footnotesize
\begin{tabular}{l r r}
\toprule
substrate & USD/QPU-sec & USD/CHSH battery \\
          &             & (4 circuits, 1024 shots) \\
\midrule
\qmirror{} (Aer + ANU)         & \textbf{0.00} & \textbf{0.00} \\
IBM Heron r2 \texttt{ibm\_fez} & 1.60--2.30    & 3.20 (observed) \\
IonQ Aria-1                    & n/a (per-shot) & 81.20 (250 sh/set) \\
IonQ Forte-1                   & n/a (per-shot) & 33.20 (100 sh/set) \\
Rigetti Cepheus 108Q           & n/a (per-shot) & 2.94  (1024 sh/set) \\
\bottomrule
\end{tabular}
\end{table}

\begin{figure}[htbp]
\centering
\begin{lstlisting}[basicstyle=\ttfamily\tiny,frame=single,framesep=4pt]
*** Figure 3 placeholder (cost-per-battery, log scale) ***

USD per 4-circuit CHSH battery (log10 axis)
  100 +                          ##  IonQ Aria-1 (81.20)
       |                         ##
   10  +              ##         ##  IonQ Forte-1 (33.20)
       |              ##         ##
    1  +    ##  ##    ##         ##  Rigetti (2.94)
       |    ##  ##                  IBM Heron (3.20)
   0.1 +
       |
   0   +##                         qmirror = 0.00 (floor)
       +-----+-----+-----+-----+-----+
        qmir  IBM   Rig  Forte  Aria
\end{lstlisting}
\caption{Per-battery cost (USD, log scale) for a 4-circuit CHSH
benchmark across \qmirror{} and four real-QPU vendors. \qmirror{}'s
cost is the floor (USD~0). Replace with \texttt{matplotlib} \texttt{semilogy} bar
chart for camera-ready.}
\label{fig:cost_bar}
\end{figure}

\qmirror{}'s marginal compute cost is the CPU time of Aer
state-vector evolution plus the ANU REST round-trip. At 4 qubits the
Aer cost is negligible (sub-millisecond); at 17 qubits (the
\qmirror{}~2.0 surface code toy) it is approximately 30 seconds
total. ANU REST cost is nonzero in latency (approximately
100--500~ms per request) but zero in USD.

\subsection{One-time calibration cost}

Total calibration spend across the closure cycle was
\textbf{USD 41.34}:

\begin{table}[htbp]
\centering
\caption{One-time calibration line items.}
\label{tab:calibration_lineitems}
\footnotesize
\begin{tabular}{l r}
\toprule
line item & USD \\
\midrule
Rigetti Cepheus 108Q (cond.8 $\beta$)        & 2.94  \\
IonQ Forte-1 (cond.8 $\beta$)                & 33.20 \\
IBM Heron r2 \texttt{ibm\_fez} (cond.3)      & 3.20  \\
ANU QRNG                                     & 0.00 (free tier) \\
Aer simulation                               & 0.00 (CPU) \\
\midrule
\textbf{total}                               & \textbf{39.34} \\
\bottomrule
\end{tabular}
\end{table}

Note: the IonQ Aria-1 reference $S = 2.808$ from
\texttt{nexus\_chsh\_bell\_2026\_05\_02} (USD~81.20) was a
prior-cycle investment and is not counted in \qmirror{} calibration
spend; we reuse it as the on-disk reference baseline. The full
disclosed total of USD~41.34 in the abstract includes \$2.00 of
exploratory Braket cost folded in for completeness; the line-item
table totals USD~39.34 of CHSH-attributable spend.

\subsection{Permanent zero-cost operation}

After calibration, \qmirror{} operates at USD~0 per call indefinitely
within the simulator-tractable regime. Quarterly IonQ anchor
refreshes (approximately USD~80 per quarter for a single Bell-pair
re-baseline) are recommended in the future-work roadmap
(Section~\ref{sec:future}) as drift estimation, not as ongoing
operating cost.
